% HD-R16. GPT-6 Astra Pro /root; inherited effort not exposed.
% 21 November 2026. Actual remaining PSL2(7) tensor sources.
\section{HD-R16: the final multiplicity-three source}
\label{sec:rank-three-tensor-final-budget}

Every period-rank-nine natural sixfold source in the
multiplicity-three alternative of HD-S45 has an already
supplied ordinary Hermitian determinant. In particular that
alternative supplies none of the missing components in HD-R15.
We prove the remaining case of HD-I22 through its actual
coefficient, quotient and original-corner maps.

Thus retain
\[
 \begin{gathered}
 K=\mathbb Q(\sqrt{-7}),\quad G=S_s\times\Gamma,
       \quad V=W\otimes_KM,\\
 W=\operatorname{Aug}(K^s),\quad r=s-1\ge9,\quad\dim_KM=3,\\
 \Gamma=L\ \hbox{or}\ \{\pm I\}\times L,
       \quad L\simeq PSL_2(7).
 \end{gathered}                                           \tag{R16.1}
\]
The matrices of \(L\), their positive Hermitian form, and
their actual identification over \(K\) are those of HD-I22.
The original sphere-cover inertia is
\((\pi_j,\gamma_j)\), \(1\le j\le b\). It generates the
displayed product and has product one. Write
\(c_j=\operatorname{rank}(P_{\pi_j}\otimes\gamma_j-I)\).
The sixfold Euler formula and period rank give
\[
              \sum_jc_j=6r+6,\qquad b\ge12.              \tag{R16.2}
\]
Every true coefficient inertia has \(c_j\ge3\), as
already proved for the HD-S45 residual source.

\paragraph{1. Exact eigenvalues and the permutation-cycle formula.}
The possible orders in \(L\) are \(1,2,3,4,7\).
This also follows directly from its actual isomorphism with
\(SL_2(\mathbb F_7)/\{\pm I\}\): a lift with repeated
eigenvalue is scalar or projectively unipotent of order seven;
distinct eigenvalues in \(\mathbb F_7\) give projective
order dividing three; and conjugate eigenvalues in
\(\mathbb F_{49}\) have norm one, giving projective order
dividing four. The exact finite matrix enumeration of HD-I22
independently verifies this list in the displayed model.

All these actual three-dimensional matrices have determinant
one. Since \(i\) and \(\omega\) are absent from \(K\),
their multiplicities in a \(K\)-characteristic polynomial
are paired with those of their conjugates. Thus the spectra
of the nonidentity order-two, order-three and order-four
elements are respectively
\[
 (1,-1,-1),\qquad (1,\omega,\omega^2),\qquad(1,i,-i).
                                                               \tag{R16.3}
\]
For order seven the irreducible cubic factors over \(K\)
give \((\zeta_7,\zeta_7^2,\zeta_7^4)\) or its conjugate,
as proved through the literal field embedding in HD-I22.
Multiplication by \(-I\) gives all remaining spectra.
In particular every actual codimension-one element of
\(\Gamma\) has determinant \(-1\), and every
codimension-two element has determinant \(+1\).

For \(\gamma\in\Gamma\) put
\(f_\ell(\gamma)=\dim\ker(\gamma^\ell-I)\).
On an \(\ell\)-cycle of a permutation, a fixed vector
of \(P_\pi\otimes\gamma\) on \(K^s\otimes M\) is
determined by \(u\in\ker(\gamma^\ell-I)\), through
the successive coordinates \(u,\gamma u,\ldots\).
The actual decomposition
\(K^s=K\boldsymbol1\oplus W\) is given by the diagonal
map and the average \(s^{-1}\sum\). Consequently
\[
 \dim\operatorname{Fix}(P_\pi\otimes\gamma\mid W\otimes M)
       =\sum_{\text{cycles of }\pi}f_\ell(\gamma)-f_1(\gamma),
 \quad c=3r-\sum_{\text{cycles}}f_\ell+f_1.             \tag{R16.4}
\]
This is an identity of fixed spaces of the original matrices.

The spectra give the following complete bounds, valid for
every positive cycle length. A sign in the first column
means multiplication of the stated \(L\)-element by that
scalar. The last column follows from (R16.4) and
\(\sum\ell=s=r+1\).
\[
\begin{array}{c|c|c|c}
 \gamma&f_1&\text{upper bound for }f_\ell&c\text{ lower bound}\\\hline
 -I&0&3\ell/2&(3r-3)/2\\
 -\text{order }2&2&2\ell&r\\
 +\text{order }2&1&3\ell/2&(3r-1)/2\\
 +\text{order }3\text{ or }4&1&\ell&2r\\
 -\text{order }3&0&\ell/2&(5r-1)/2\\
 -\text{order }4&0&3\ell/4&(9r-3)/4\\
 +\text{order }7&0&3\ell/7&(18r-3)/7\\
 -\text{order }7&0&3\ell/14&(39r-3)/14
\end{array}                                                   \tag{R16.5}
\]
For example, the negative order-four element has
\(f_\ell=0,1,3\) according as \(\ell\) is odd,
is twice an odd number, or is divisible by four. The
negative order-three element has value three at multiples
of six, one at even lengths not divisible by three, and
zero at odd lengths. The seventh-root cases have value
three exactly at multiples of seven or fourteen. These
descriptions verify the bounds at every length, including
the endpoint lengths. All nonidentity \(\gamma\) have
\(c\ge r\). A noncentral element that is not a reflection
has \(c\ge(3r-1)/2\); every element of projective order
at least three has \(c\ge2r\), since \(r\ge9\).

\paragraph{2. At most five nontrivial entries, and none scalar.}
Let \(k\) count the nonidentity \(\gamma_j\), and let
\(z\) count occurrences of \(-I\). The projective tuple
generates \(PSL_2(7)\). At least three entries are
projectively nonidentity: a product-one tuple with at most
two such entries generates a cyclic projective group.
Thus \(k-z\ge3\). The bounds above give
\[
 6r+6\ge kr+3(b-k)\ge k(r-3)+36,
 \qquad k\le5.                                           \tag{R16.6}
\]
Indeed \(k\ge6\) would give a lower bound
\(6r+18>6r+6\), because \(r-3>0\).

There are no central nonidentity entries. For \(z\ge2\),
(R16.6) and \(k-z\ge3\) force \(k=5,z=2\).
The two scalar entries and the other three then give
\[
 \sum_jc_j\ge 2(3r-3)/2+3r+3(12-5)=6r+18,
                                                               \tag{R16.7}
\]
a contradiction. For \(z=1,k=4\), the three projective
nonidentity entries have product one and generate
\(PSL_2(7)\). At most one of them is an involution.
In fact two involutions \(a,b\) with the third equal
to \((ab)^{-1}\) generate a dihedral group: every word
reduces to \((ab)^j\) or \(a(ab)^j\). Such a group
has a cyclic subgroup of index at most two, whereas
\(PSL_2(7)\) is nonabelian perfect. Hence at least two
entries have projective order at least three, giving
\[
 \sum_jc_j\ge(3r-3)/2+r+4r+3(12-4)
                 =(13r+45)/2>6r+6.                       \tag{R16.8}
\]
For \(z=1,k=5\), the four noncentral entries cannot
all be reflections: their determinant product would be
\(+1\), and multiplication by the determinant of \(-I\)
would contradict the literal product-one relation.
At least one is a nonreflection, so
\[
 \sum_jc_j\ge(3r-3)/2+(3r-1)/2+3r+3(12-5)
                      =6r+19>6r+6.                       \tag{R16.9}
\]
These exhaust the possibilities. Therefore \(z=0\).

\paragraph{3. The actual base summand has dimension at most one.}
Let \(d_0\) be the \(K\)-dimension of the parabolic
cohomology of \(M\) on the original sphere with monodromies
\(\gamma_j\). The connected actual \(\Gamma\)-quotient
cover has this coefficient system, which is absolutely
irreducible and nontrivial. Its two invariant spaces vanish,
so Euler gives the integer
\[
 d_0=-6+\sum_j\operatorname{rank}(\gamma_j-I)\ge0.       \tag{R16.10}
\]
Since all \(k\) nonidentity entries are noncentral, let
\(a,c,d\) count those with moving ranks one, two, three
on \(M\), respectively. Then
\[
 a+c+d=k\in\{3,4,5\},\qquad
              d_0=k+c+2d-6.                             \tag{R16.11}
\]
The codimension-three rows of (R16.5) all have lower
bound at least \((9r-3)/4\). Thus
\[
 \sum_{\gamma_j\ne1}c_j
   \ge r\bigl(k+c/2+5d/4\bigr)-c/2-3d/4.              \tag{R16.12}
\]
We exclude \(d_0\ge2\) by an exact finite integer
inequality. It would give \(k+c+2d\ge8\). For \(k=3\)
this requires \(d\ge2\), and
\(k+c/2+5d/4\ge6\). For \(k=4\) the same expression
is at least \(6+d/4\); for \(k=5\) it is at least
\(13/2+d/4\). Also
\(c/2+3d/4\le3k/4\le15/4\). Combining these inequalities
with the other \(b-k\ge7\) inertia entries yields
\[
 \sum_jc_j\ge6r-15/4+21
                       =6r+69/4>6r+6,                   \tag{R16.13}
\]
contrary to (R16.2). We have proved
\[
                              d_0\in\{0,1\}.            \tag{R16.14}
\]

\paragraph{4. Actual quotient maps and the final period calculation.}
We recall the precise curves to keep the original factor
through this calculation. For the original connected
\(G\)-cover \(C\) put
\[
 \begin{gathered}
 L_0=S_{s-1},\quad X=C/(L_0\times1),\quad Z=C/(S_s\times1),\\
 Y=C/(L_0\times\Gamma),\quad
 \alpha:C\to X,\quad f:X\to Z,\quad\deg f=s.
 \end{gathered}                                           \tag{R16.15}
\]
The map \(X\to Y\) is an actual connected
\(\Gamma\)-cover. The diagonal and averaging maps on
\(K^s\) split its rank-three parabolic cohomology as
\[
 H^1_{\rm par}(Y,M)
   =H^1_{\rm par}(\mathbb P^1,M)
        \mathbin{\perp}H^1_{\rm par}(\mathbb P^1,W\otimes M).
                                                               \tag{R16.16}
\]
The first summand has dimension \(d_0\), the second
dimension six, over \(K\); the first pairing is scaled
by \(s\). At a ramified sheet cycle both maps use exactly
\(\gamma^\ell\), as in (R16.4), so these are the
actual middle-extension and paired Hodge maps of HD-S45.

First carry the originally prescribed primitive vector
by the established integer-cleared Fourier matrix unit
to \(w_0\otimes u\), where \(w_0\) spans the
\(L_0\)-fixed augmentation line. For the primitive
\(\Gamma\)-corner \(e_u\), the complementary projector is
\[
 \pi=e_u(1-s^{-1}f^*f_*),\qquad
 |L_0|^{-1}\alpha^*\pi\alpha_*
       =(\operatorname{av}_{L_0}-\operatorname{av}_{S_s})e_u.
                                                               \tag{R16.17}
\]
These are actual rational Jacobian operators. Their
integer-cleared images and the quotient pullback give
the original sixfold through the isogenies of HD-S45 and
HD-I22, with multipliers \(|L_0|\) and
\(M_0^2h(v,v)/h(w,w)\) for the change of primitive line.
Every such positive rational multiplier has sixth power
equal to the norm of its cube in \(K\), so the prescribed
ordinary determinant class is retained.

Over a finite parameter cover fixing the finite group
labels, all these maps are flat maps of the paired local
systems and their Hodge filtrations. In particular the
period differential of the sixfold summand is obtained
by applying the constant projector (R16.17) to that of
the rank-three system on \(Y\). Its rank cannot exceed
the rank of the latter period differential.

For \(d_0=0\), the connected primitive image
\(e_uJZ\) has zero first cohomology and is the zero
abelian variety. The \(\Gamma\)-equivariant norm satisfies
\(f_*e_u=e_uf_*\), so its image on this corner is zero.
Consequently \(f^*f_*e_u=0\) and the actual rational
operator (R16.17) satisfies \(\pi=e_u\) on \(JX\).
Thus (R16.16)--(R16.17) identify the whole rank-three
quotient factor with the transported sixfold.
Its Euler equation is \(6=3(2g(Y)-2)+\sum e_j\).
The genus-one and genus-two marked-base bounds are six
and three. In genus zero there are at most twelve genuine
coefficient branch values. Period rank nine therefore
forces \(Y=\mathbb P^1\), twelve such values, and
\(e_j=1\) for all of them. The actual quotient group
is generated by reflections. HD-G27, equivalently the
rank-three endpoint of HD-S40, now applies through this
literal source and gives ordinary determinant \([-1]\).
This is an already supplied component; the subgroup
\(\Gamma=L\), which has no reflections, cannot occur
in this period-rank-nine endpoint.

For \(d_0=1\), the rank-three quotient has parabolic
dimension seven. If \(g=g(Y)\), and \(t\) counts its
genuine coefficient branch values with moving ranks
\(e_j\ge1\), its exact Euler formula is
\[
                    7=3(2g-2)+\sum_{j=1}^t e_j.         \tag{R16.18}
\]
For \(g=1\) its marked-base parameter dimension is at
most seven. For \(g=2\) it is at most four, and larger
genus is impossible. In genus zero the sum is thirteen.
A period rank of at least nine would require \(t\ge12\).
The only possibilities are thirteen rank-one inertias,
or eleven rank-one inertias and one rank-two inertia.
In both cases their determinant product is \(-1\),
by the complete actual spectra (R16.3)--(R16.5).
The sphere monodromy relation requires determinant product
\(+1\), a contradiction. Thus \(t\le11\), and the
rank-three quotient family has period rank at most eight.
Finite monodromy labels give finite discrete cover data
over each marked base; omitted coefficient-trivial branch
values add no parameters to this quotient variation.
Equation (R16.17) gives the same bound for the prescribed
sixfold summand.

The two values in (R16.14) therefore give, respectively,
the already supplied determinant or period rank at most
eight. Combining this exact conclusion with HD-I22 and
the earlier HD-S45 cases closes the complete
multiplicity-three alternative for the missing sixfold
components. No existence or dominance assertion has been
inferred from any parameter dimension.
